\section{舵机执行通路异步化的可行性分析}\label{sec:async-write}

\subsection{半双工 TTL 总线的物理约束}\label{sec:async-write-bus}

机械臂六个 STS-3215 舵机走的是 Feetech SCServo 协议，物理层是
单根 TTL 半双工串行线（/dev/ttyACM0）。
``半双工''意味着同一时刻总线上只能存在一个方向的传输：
要么主机向某个 ID 写指令，要么某个舵机向主机回报状态。
六个舵机和主机控制板共享这一根线，所有事务必须严格串行。

这与摄像头的 USB 链路有本质区别。每台 USB 摄像头独占一条 USB 接口，
后台异步读取线程不停读取不会影响其他设备；
而舵机只要后台线程发起任何写入，就会和主线程的同步读取
在物理总线上撞车，造成字节流交错、回包校验失败、整帧数据废掉。

\subsection{异步写为什么不可行}\label{sec:async-write-why-not}

把同步写入改为后台队列异步执行，至少要面对四条独立的约束，
任意一条都足以让该方案失败：

约束 A——总线竞争：异步写线程和主线程的同步读取
不可避免地会同时尝试占用总线。即使引入互斥锁让两者轮流持锁，
锁竞争和上下文切换会把异步带来的并发收益完全吃掉，
退化成纯串行加额外开销，等于没异步。

约束 B——安全限幅链断裂：在
so101\_follower.py 的 send\_action 函数中，
当用户配置了 max\_relative\_target 时存在如下读改写链：
系统先将策略动作解析为目标位置；若启用了最大相对位移限制，
则必须先读取当前舵机位置，再根据当前位置对目标位置做安全限幅，
最后把限幅后的目标位置同步写入舵机总线。
这条链严格要求``读到的当前位置''与``写出的目标位置''属于同一时刻，
限幅基准才有意义。一旦写操作异步化，下一帧的同步读取
可能读到上一帧目标尚未执行完的中间位置，
安全目标位置检查的限幅基准变得不确定，
最坏情况下机械臂可能在限幅阈值附近发生连续误判。

约束 C——错误暴露延迟：当前同步写入接口在
通信失败时立即抛出连接错误，调用方
（机器人动作发送接口）能在同一帧内检测到``指令没送达''，
可以停机、重试或上报。如果改为异步：主线程已经返回成功，
policy 继续向前推理，几十毫秒后后台线程才发现包掉了，
此时机械臂可能已经卡在某个半执行姿态，policy 还在按
``执行成功''的假设输出下一条指令。在机器人安全场景里，
这种错误反馈的滞后是不可接受的。

约束 D——收益不在瓶颈上：参考第三章的耗时分布，
同步写入一帧约 1--3 ms，而 ACT 推理一次约 2000 ms。
即便假定异步写能完美隐藏掉这 1--3 ms 的阻塞，
对总循环时间的削减也不到 0.1\%。在当前推理主导的负载下，
任何 ROI 计算都会得到接近零的结果。

\subsection{小结}\label{sec:async-write-conclusion}

四条约束分别来自硬件协议（A）、应用层不变量（B）、
安全工程实践（C）和工程 ROI（D），任何一条都构成独立的否决理由。
更关键的是，约束 A 和 B 是结构性的：只要硬件保持半双工 TTL 总线，
只要动作发送接口还要做读改写式的安全限幅，
异步写在本平台上都不存在落地空间。LeRobot 上游代码中既没有
异步写入占位方法、也没有任何 TODO 注释，
也是出于同样的判断。因此本论文不再把舵机异步写视为候选优化方向。
